\documentclass[11pt,a4paper]{article}

\usepackage[margin=1in]{geometry}
\usepackage{hyperref}
\usepackage{enumitem}
\usepackage{mdframed}
\usepackage{booktabs}
\usepackage{tabularx}

\hypersetup{colorlinks=true, linkcolor=blue!70!black, urlcolor=blue!70!black}

\title{ERC Indoor Laser Focus\\[0.3em]
\large What we implement next and what we deliberately postpone}
\author{Project Working Draft}
\date{\today}

\begin{document}

\maketitle
\begin{center}
\fbox{\parbox{0.9\textwidth}{
\textbf{HISTORICAL DOCUMENT} \\[4pt]
Scope document from March 2026. The scope decisions here were either implemented,
dropped, or folded into later runtime work. SuperGlue appears in this file but
was never part of the final indoor backbone; CosPlace alone proved sufficient. \\[4pt]
\textit{For the current system, see: \texttt{CLAUDE.md}}
}}
\end{center}
\vspace{1em}

\section{One-Sentence Project Definition}

We are building a \textbf{known-corridor indoor navigation system} for ERC using:

\begin{itemize}[nosep]
  \item \textbf{CosPlace + SuperGlue + graph} for perception and planning,
  \item a \textbf{motion-prior layer} built from orientation / IMU support,
  \item a \textbf{validated online local controller} for nearby graph-node movement,
  \item a separate \textbf{safety layer},
  \item a separate \textbf{checkpoint verification and recovery layer}.
\end{itemize}


\section{What We Are Using Right Now}

\subsection*{Core baseline}

\begin{itemize}[nosep]
  \item \texttt{baseline.py} is the perception/planning baseline.
  \item It is responsible for descriptor extraction, retrieval, graph construction, and query-time localization support.
  \item It is the right starting point because our corridor is known and repeated.
\end{itemize}

\subsection*{Robot interface}

\begin{itemize}[nosep]
  \item \texttt{src/earthrover\_interface.py} gives us the live robot interface.
  \item The front camera frame is still the main localization signal.
  \item The SDK also exposes orientation, speed, and IMU arrays (\texttt{accels}, \texttt{gyros}, \texttt{mags}).
  \item For indoor ERC, those motion signals should be used as \textbf{supporting priors} for heading consistency, node-belief smoothing, and recovery.
  \item GPS is available in the SDK, but it is not part of the indoor backbone.
\end{itemize}

\subsection*{Depth safety}

\begin{itemize}[nosep]
  \item \texttt{src/depth\_estimator.py} and \texttt{src/depth\_safety.py} are optional safety tools.
  \item They are useful after the baseline localization/planning loop is working.
\end{itemize}

\subsection*{MBRA / LogoNav code}

\begin{itemize}[nosep]
  \item \texttt{mbra\_repo/} is a research reference, not yet our deployed indoor controller.
  \item MBRA is the relabeling expert side of that project.
  \item LogoNav is the deployed-policy side, but in this repo it is GPS/pose-conditioned, so it is not a drop-in indoor image-goal controller.
\end{itemize}


\section{What We Are \underline{Actually} Doing First}

\begin{mdframed}
\textbf{Immediate focus:} stop thinking about the whole final system at once.
The next job is to make the corridor baseline real on our actual data.
\end{mdframed}

\subsection*{Step 1: Collect corridor data}

\begin{itemize}[nosep]
  \item Run the robot through the real corridor.
  \item Save front-camera images in order.
  \item Save orientation, speed, IMU samples, timestamps, and any available command / teleoperation metadata per step.
  \item Collect multiple runs of the same route.
\end{itemize}

\textbf{Deliverable:} one clean reference dataset from the real ERC corridor with both images and motion-state logs.

\subsection*{Step 2: Build the corridor database}

\begin{itemize}[nosep]
  \item Run \texttt{baseline.py build-db}.
  \item Build descriptors.
  \item Build the place graph.
  \item Build the action graph if route metadata is available.
\end{itemize}

\textbf{Deliverable:} database artifacts that can be queried repeatedly.

\subsection*{Step 3: Validate localization}

\begin{itemize}[nosep]
  \item Run \texttt{baseline.py query} on real corridor images.
  \item Check whether retrieval gives the correct corridor nodes.
  \item Identify ambiguous segments, weak landmarks, and failure cases.
  \item Tune graph density and image preprocessing.
  \item Add heading-consistency checks so retrieval cannot jump to impossible nodes too easily.
\end{itemize}

\textbf{Deliverable:} a corridor graph and localization pipeline that is stable enough to trust, including temporal and heading-aware filtering.

\subsection*{Step 4: Decide the online local controller}

\begin{itemize}[nosep]
  \item Decide what actually runs online for nearby graph-node movement.
  \item MBRA / LogoNav is a reference here, but should not be assumed correct by default.
  \item The decision must be made from real corridor testing, not from paper naming.
  \item The chosen controller should use heading / yaw-rate feedback for smoothing if that improves stability.
\end{itemize}

\textbf{Deliverable:} one local controller candidate that we can actually test online.

\subsection*{Step 5: Integrate runtime loop}

\begin{itemize}[nosep]
  \item camera frame,
  \item localization in graph,
  \item graph path / next subgoal,
  \item online local controller,
  \item safety override,
  \item checkpoint verification.
\end{itemize}

\textbf{Deliverable:} one full indoor loop for a single checkpoint.


\section{What We Should \underline{Not} Split Attention On Right Now}

\begin{itemize}[nosep]
  \item training a brand new large model,
  \item replacing the graph-localization baseline before it is tested,
  \item making SLAM the backbone before the current baseline is evaluated on the real corridor,
  \item polishing multi-checkpoint logic before single-checkpoint execution works,
  \item deep safety tuning before localization and control are running.
\end{itemize}


\section{Sensors and Signals: What Matters Now}

\begin{tabularx}{\textwidth}{p{0.22\textwidth} p{0.18\textwidth} X}
\toprule
\textbf{Signal} & \textbf{Priority} & \textbf{How we use it now} \\
\midrule
Front camera & Highest & Main signal for localization, graph matching, checkpoint verification, and local control input \\
Orientation / heading & High & Motion prior for heading consistency, candidate-node gating, subgoal alignment, and recovery behavior \\
IMU arrays & High-medium & Short-term yaw / motion support for temporal filtering, controller damping, and uncertainty handling \\
GPS & Low for indoor & Mostly irrelevant for the main indoor corridor task \\
Wheel encoders / RPMs & Validate first & Use only if live tests prove they are populated and stable; do not assume odometry quality without verification \\
Depth estimate & Medium-later & Safety support after baseline localization/control is working \\
\bottomrule
\end{tabularx}

\textbf{Practical answer:} the first working indoor baseline should be
\textbf{vision-primary with IMU / heading support}. It should not be vision-only,
and it should not assume encoder-based odometry until the robot proves that signal is real.


\section{The Main Open Question}

The biggest unresolved implementation question is:

\begin{center}
\fbox{\parbox{0.85\textwidth}{
\textbf{What is our actual online short-horizon controller for moving between nearby graph nodes?}
}}
\end{center}

Everything else has a cleaner answer already:

\begin{itemize}[nosep]
  \item perception/planning baseline = \texttt{baseline.py}
  \item robot interface = \texttt{src/earthrover\_interface.py}
  \item motion prior = orientation / IMU support with temporal filtering
  \item safety support = depth and simple runtime veto logic
\end{itemize}


\section{Immediate Milestones}

\subsection*{Milestone 1}

\textbf{Collect corridor data and build the reference database.}
\textbf{At the same time, verify orientation / IMU / RPM availability on the real robot.}

\subsection*{Milestone 2}

\textbf{Prove that query-time localization works reliably on repeated corridor runs.}
\textbf{Add temporal and heading-aware filtering to the node belief.}

\subsection*{Milestone 3}

\textbf{Choose and validate one online local controller on nearby graph-node subgoals.}

\subsection*{Milestone 4}

\textbf{Run one single-checkpoint indoor loop with localization + planning + controller + safety.}

\subsection*{Milestone 5}

\textbf{Add checkpoint verification and multi-checkpoint execution.}


\section{What Success Looks Like Right Now}

Success in the next phase is \textbf{not}:

\begin{itemize}[nosep]
  \item a polished final competition stack,
  \item a fully trained new model,
  \item a fancy architecture diagram.
\end{itemize}

Success in the next phase \textbf{is}:

\begin{itemize}[nosep]
  \item we can collect corridor data cleanly,
  \item we can log the motion signals we actually trust,
  \item we can build a graph database from it,
  \item we can localize repeated query frames correctly,
  \item we can stabilize localization using heading / IMU support,
  \item we can pick the next nearby graph node,
  \item we have one real candidate for the online local controller.
\end{itemize}


\section{Bottom Line}

\begin{mdframed}
\begin{itemize}[nosep]
  \item The baseline is ready enough to start real work.
  \item The next step is \textbf{data collection on the actual corridor}.
  \item The project should be \textbf{vision-primary with IMU / heading support}.
  \item We must verify which onboard motion signals are actually trustworthy.
  \item The biggest open engineering decision is the \textbf{online local controller}.
  \item Everything else should stay focused on making the known-corridor baseline work end-to-end.
\end{itemize}
\end{mdframed}

\end{document}
